\documentclass[10pt,letterpaper]{article}

\usepackage[
	letterpaper,
	left=0.48in,
	right=0.48in,
	top=0.40in,
	bottom=0.36in
]{geometry}
\usepackage[T1]{fontenc}
\usepackage[utf8]{inputenc}
\usepackage{newtxtext}
\usepackage{enumitem}
\usepackage{tabularx}
\usepackage{titlesec}
\usepackage[hidelinks]{hyperref}
\usepackage{fontawesome5}
\usepackage{xurl}
\usepackage{microtype}

\pagestyle{empty}
\raggedbottom
\setlength{\parindent}{0pt}
\setlength{\parskip}{0pt}
\setlength{\tabcolsep}{0pt}
\setlength{\emergencystretch}{2em}
\urlstyle{same}

\IfFileExists{glyphtounicode.tex}{%
	\input{glyphtounicode}
	\pdfgentounicode=1
}{}

\titleformat{\section}
	{\bfseries\large}{}{0pt}{}
	[\vspace{1.2pt}\titlerule]
\titlespacing*{\section}{0pt}{5.0pt}{3.0pt}

\newcommand{\resumeEntry}[4]{%
	\begin{tabularx}{\textwidth}{@{}Xr@{}}
		\textbf{#1} & #2 \\
		\textit{#3} & \textit{#4}
	\end{tabularx}%
}

\newcommand{\projectEntry}[2]{%
	\begin{tabularx}{\textwidth}{@{}Xr@{}}
		\textbf{#1} & \textit{#2}
	\end{tabularx}%
}

\newenvironment{resumeItems}{%
	\begin{itemize}[
		label=\textbullet,
		leftmargin=14pt,
		labelsep=5pt,
		itemsep=0.5pt,
		topsep=1.0pt,
		parsep=0pt,
		partopsep=0pt
	]
}{\end{itemize}}

\newcommand{\contactsep}{\qquad}

\begin{document}
\normalsize

\begin{center}
	{\LARGE\textbf{Nahid Hasan Sagar}}\\[2.2pt]
	{\footnotesize
		\faMapMarker*\, Bloomington, Indiana
		\contactsep
		\href{tel:+19309041643}{\faPhone*\, +1 (930) 904-1643}
		\contactsep
		\href{mailto:nahidhasansagar@gmail.com}{\faEnvelope\, nahidhasansagar@gmail.com}
	}\\[-0.2pt]
	{\footnotesize
		\href{https://www.linkedin.com/in/nahidhasansagar/}{\faLinkedinIn\, linkedin.com/in/nahidhasansagar}
		\contactsep
		\href{https://github.com/Sagar-DU}{\faGithub\, github.com/Sagar-DU}
	}
\end{center}

\vspace{-5pt}

\section*{PROFESSIONAL SUMMARY}

M.S. candidate in Statistical Science at Indiana University Bloomington with a B.S. in Physics and experience developing quantitative research software in Python, R, and C++. Applies Bayesian modeling, statistical learning, time-series analysis, and numerical methods to data-intensive problems, with a focus on rigorous validation, reproducibility, and computational performance.

\section*{EDUCATION}

\resumeEntry
	{Indiana University, College of Arts and Sciences}
	{Bloomington, IN}
	{Master of Science in Statistical Science}
	{Expected May 2027}

\vspace{0.4pt}
\textit{Thesis: Time-Series Analysis Using Bayesian Additive Regression Trees (BART)}

\vspace{0.5pt}
\textbf{Relevant Coursework:} Introduction to Statistical Computing; Applied Statistical Computing; Fundamentals of Statistics I--II; Applied Linear Models I--II; Statistical Consulting; Research in Statistics; Introduction to Analysis I

\vspace{2.2pt}
\resumeEntry
	{University of Dhaka}
	{Dhaka, Bangladesh}
	{Bachelor of Science in Physics}
	{August 2024}

\section*{RESEARCH EXPERIENCE}

\resumeEntry
	{Research Assistant, StatMIND Lab}
	{2025 -- Present}
	{Indiana University Bloomington}
	{Bloomington, IN}

\begin{resumeItems}
	\item Develop Bayesian brain-mapping algorithms for high-dimensional neuroimaging data, translating statistical research methods into efficient C++ implementations integrated with R through Rcpp and RcppArmadillo.
	\item Reproduce results from the original R package and validate the new implementation with unit tests, benchmark comparisons, and systematic numerical checks.
	\item Achieve a 4--5$\times$ speedup in core computations, enabling faster iteration on large-scale datasets.
	\item Build reproducible workflows to compare runtime, scalability, and computational performance across implementations.
\end{resumeItems}

\section*{SELECTED RESEARCH \& PROJECTS}

\projectEntry
	{M.S. Thesis: Time-Series Analysis Using BART}
	{R, Bayesian Modeling}
\begin{resumeItems}
	\item Conduct research on Bayesian Additive Regression Trees for time-series analysis under the supervision of Dr. Matthew T. Pratola, evaluating model behavior for structured temporal data.
\end{resumeItems}

\vspace{0.7pt}
\projectEntry
	{Ordinary Least Squares from Scratch}
	{Python, NumPy, Pandas, Matplotlib}
\begin{resumeItems}
	\item Implement ordinary least squares regression from first principles using matrix-based estimation, with transparent computation and data-analysis visualizations.
\end{resumeItems}

\section*{TECHNICAL SKILLS}

\textbf{Programming:} Python, R, C++, C, MATLAB, Bash\\
\textbf{Quantitative Methods:} Bayesian inference, BART, MCMC, time-series analysis, linear and logistic regression, mixed-effects models, bootstrap methods, numerical optimization, model selection, cross-validation, classification metrics, predictive model evaluation\\
\textbf{Research Engineering:} Algorithm benchmarking, code profiling, unit testing, reproducible workflows, Git, GitHub, LaTeX, RStudio, VS Code

\section*{LEADERSHIP \& HONORS}

\begin{resumeItems}
	\item Attended the 2026 Velocity Conference at UC Berkeley as an Indiana University representative for the science community.
	\item Co-organized the Australian Mathematics Competition in Bangladesh in 2014.
	\item Taught and mentored more than 100,000 students through classroom instruction and national online learning platforms as a Senior Mathematics Instructor.
\end{resumeItems}

\end{document}
